% from ega4-18-fr.tex
\oldpage[IV-4]{179}

Et en localisant en $\mathfrak p$, il vient un isomorphisme
\[
  (\hat{\Omega}^1_{A/k})_{\mathfrak p}\xrightarrow{\sim}(\hat{\Omega}^1_{A_1/k})_{\mathfrak p_1}.
\]
Ceci montre que la condition $b)$ pour l'anneau $A$ et l'idéal $\mathfrak p$ est équivalente à la condition $b)$ pour l'anneau $A_1$ et l'idéal $\mathfrak p_1$. D'autre part, avec les notations de $a)$, on a $A'_1=A_1\widehat{\otimes}_k k'=A'/\mathfrak JA'$ à isomorphisme près ((7.5.7.1), où l'hypothèse sur le corps résiduel de $B$ est superflue); si $\mathfrak p'$ est un idéal premier de $A'$ au-dessus de $\mathfrak p$, tout élément de $\mathfrak JA'$ annule un élément de $A'-\mathfrak p'$, donc $\mathfrak JA'\subset\mathfrak p'$, et si l'on pose $\mathfrak p'_1=\mathfrak p'/\mathfrak JA'$, $\mathfrak p'_1$ est au-dessus de $\mathfrak p_1$ et $(A'_1)_{\mathfrak p'_1}$ s'identifie canoniquement à $A'_{\mathfrak p'}$; ceci montre donc que la condition $a)$ (resp. $a')$) pour l'anneau $A$ et l'idéal $\mathfrak p$ est équivalente à la condition $a)$ (resp. $a')$) pour l'anneau $A_1$ et l'idéal $\mathfrak p_1$. Or, tous les idéaux premiers minimaux de $A$ contenus dans $\mathfrak p$ contiennent $\mathfrak J$ puisque $\Spec(A_1)\to\Spec(A)$ est un isomorphisme local au point $\mathfrak p_1$; d'autre part les idéaux de $\operatorname{Ass}_A(A/\mathfrak J)$ sont les idéaux de $\operatorname{Ass}(A)$ qui sont contenus dans $\mathfrak p$ (Bourbaki, \emph{Alg. comm.}, chap. IV, \textsection~1, n\textsuperscript{o}~2, prop. 6); donc les idéaux premiers minimaux de $A_1$ sont tous contenus dans $\mathfrak p_1$, et on a par suite $\dim(A_1)=n$. Il suffit alors d'appliquer à $A_1$ et à $\mathfrak p_1$ ce qui a été prouvé plus haut.

\begin{remarks}[18.11.11]
\label{IV.18.11.11}
\begin{enumerate}
  \item[(i)] L'équivalence des conditions $d)$ et $b)$ dans (18.11.10) n'est plus valable lorsqu'on ne suppose plus que $[k:k^p]<+\infty$. En effet, dans l'exemple de (0, 22.7.7, (ii)), l'anneau $B=A/\mathfrak q$ est intègre et de dimension 1; d'autre part, la suite
  \[
    (\mathfrak q/\mathfrak q^2)\otimes_A L\xrightarrow{j}\hat{\Omega}^1_{A/k}\otimes_A L\longrightarrow\hat{\Omega}^1_{B/k}\otimes_B L\longrightarrow0
  \]
  est exacte (même démonstration que dans (18.11.3)), et on sait que $j$ n'est pas injectif, donc $j=0$ puisque $(\mathfrak q/\mathfrak q^2)\otimes_A L$ est de rang 1. On en conclut que $\operatorname{rg}_L(\hat{\Omega}^1_{B/k}\otimes_B L)=2$. Comme $B_{\mathfrak p}$ est une $k$-algèbre géométriquement régulière, on voit qu'ici la condition $d)$ n'entraîne pas $b)$.
  \item[(ii)] Les notations étant celles de (18.11.10), supposons que $n=\dim(A)$ et soit $(x_i)_{1\leq i\leq n}$ un système de paramètres de $A$ n'appartenant pas à $\mathfrak p$; on en déduit un $k$-homomorphisme local $u:B\to A$ tel que $u(T_i)=x_i$ pour $1\leq i\leq n$, faisant de $A$ une $B$-algèbre finie. Pour que le morphisme correspondant $\Spec(A)\to\Spec(B)$ soit étale au point $\mathfrak p$, il faut et il suffit que les images des $d_Ax_i$ dans $(\hat{\Omega}^1_{A/k})_{\mathfrak p}$ forment un système de générateurs de cet $A_{\mathfrak p}$-module. En effet, on a vu dans la démonstration de (18.11.5) que si le morphisme $\Spec(A)\to\Spec(B)$ est non ramifié au point $\mathfrak p$, l'homomorphisme canonique $(\hat{\Omega}^1_{B/k}\otimes_B A)_{\mathfrak p}\to(\hat{\Omega}^1_{A/k})_{\mathfrak p}$ est surjectif et les images des éléments $d_BT_i\otimes1$, qui engendrent $(\hat{\Omega}^1_{B/k}\otimes_B A)_{\mathfrak p}$, sont les $d_Ax_i$; d'où la nécessité de la condition. Inversement, le même raisonnement montre que si cette condition est vérifiée, le morphisme $\Spec(A)\to\Spec(B)$ est non ramifié en $\mathfrak p$, et le raisonnement de (18.11.9) prouve en fait que ce morphisme est étale en $\mathfrak p$.
\end{enumerate}
\end{remarks}

\begin{corollary}[18.11.12]
\label{IV.18.11.12}
Soient $k$ un corps d'exposant caractéristique $p$ tel que $[k:k^p]<+\infty$, $A$ une $k$-algèbre locale noethérienne complète et intègre, qui n'est pas un corps, et dont le corps résiduel est une extension finie de $k$. Il existe alors une extension finie radicielle $k'$ de $k$ telle qu'en posant $A'=A\otimes_k k'$, la $k'$-algèbre $A'_{\mathrm{red}}$ et l'idéal premier $0$ de cette algèbre vérifient les conditions équivalentes $a)$, $a')$, $b)$, $c)$ et $d)$ de (18.11.10). En particulier, si $n=\dim(A)=\dim(A')$,

